%2multibyte Version: 5.50.0.2960 CodePage: 65001
%\usetheme{Madrid}
%\AtBeginSection[]{
%\begin{frame}
%\frametitle{Indice}
%\vfill
%\tableofcontents[sectionstyle=show/hide,subsectionstyle=show/show/hide]
%\vfill
%\end{frame}
%}

\documentclass[trans]{beamer}%
\usepackage[italian]{babel}
\usepackage{eurosym}
\usepackage{amsfonts}
\usepackage{mathpazo}
\usepackage{multimedia}
\usepackage{url}
\usepackage{amsmath}
\usepackage{amssymb}
\usepackage{graphicx}%
\setcounter{MaxMatrixCols}{30}
%TCIDATA{OutputFilter=latex2.dll}
%TCIDATA{Version=5.50.0.2960}
%TCIDATA{Codepage=65001}
%TCIDATA{CSTFile=beamer.cst}
%TCIDATA{Created=Monday, May 30, 2011 12:02:18}
%TCIDATA{LastRevised=Sunday, October 11, 2026 17:01:36}
%TCIDATA{<META NAME="GraphicsSave" CONTENT="32">}
%TCIDATA{<META NAME="SaveForMode" CONTENT="1">}
%TCIDATA{BibliographyScheme=Manual}
%TCIDATA{<META NAME="DocumentShell" CONTENT="Other Documents\SW\Slides - Beamer">}
%TCIDATA{Language=American English}
%BeginMSIPreambleData
\providecommand{\U}[1]{\protect\rule{.1in}{.1in}}
%EndMSIPreambleData
\newenvironment{stepenumerate}{\begin{enumerate}[<+->]}{\end{enumerate}}
\newenvironment{stepitemize}{\begin{itemize}[<+->]}{\end{itemize} }
\newenvironment{stepenumeratewithalert}{\begin{enumerate}[<+-| alert@+>]}{\end{enumerate}}
\newenvironment{stepitemizewithalert}{\begin{itemize}[<+-| alert@+>]}{\end{itemize} }
\usetheme{Madrid}
\definecolor{marked_yellow}{rgb}{0.95,0.5,0}
\definecolor{light-blue}{rgb}{0.50,0.55,1.00}

\subtitle{Metodi Quantitativi per la Finanza}
\institute[]{{\large Dipartimento di Economia e Management}}
\titlegraphic{
% L'interpolazione PDF migliora la resa a schermo del PNG ridimensionato.
\pdfximage width 3.3cm attr {/Interpolate true}
  {E:/Didattica/MQF/Github/CORSO-MQF/graphics/logo-Unibs-esteso.png}
\pdfrefximage\pdflastximage
}
\begin{document}

\title{Presentazione}
\author[Falbo P., Pelizzari C.]{Paolo Falbo, Cristian Pelizzari}
\date{}
\maketitle%

%TCIMACRO{\TeXButton{define light-blue}{\definecolor{light-blue}{rgb}%
%{0.50,0.55,1.00}}}%
%BeginExpansion
%
%EndExpansion
%

%TCIMACRO{\TeXButton{light-blue}{\color{light-blue}} }%
%BeginExpansion
\color{light-blue}
%EndExpansion
%TCIMACRO{\TeXButton{black}{\color{black}}}%
%BeginExpansion
\color{black}%
%EndExpansion
\medskip

\section{Cosa, chi, quando, dove, perch\'{e}, come}%

%TCIMACRO{\TeXButton{BeginFrame}{\begin{frame}}}%
%BeginExpansion
\begin{frame}%
%EndExpansion
%

%TCIMACRO{\QTR{frametitle}{Cosa, chi, quando, dove, perch\'{e}, come}}%
%BeginExpansion
\frametitle{Cosa, chi, quando, dove, perch\'{e}, come}%
%EndExpansion
%

%TCIMACRO{\QTR{framesubtitle}{Cosa, chi, quando, dove}}%
%BeginExpansion
\framesubtitle{Cosa, chi, quando, dove}%
%EndExpansion
%

%TCIMACRO{\TeXButton{Transparency}{\setbeamercovered{transparent=20}}}%
%BeginExpansion
\setbeamercovered{transparent=20}%
%EndExpansion


\begin{stepitemize}
\item \textit{Metodi Quantitativi per la Finanza} (\textit{MQF})

\begin{stepitemize}
\item insegnamento del Corso di Laurea Magistrale in Finanza e Risk
Management, codice ECO0139

\item docenti: Falbo Paolo, Pelizzari Cristian
\end{stepitemize}

\item \textit{Lezioni}

\begin{stepitemize}
\item luned\`{\i} 14:30 - 16:45, laboratorio informatico 3 della sede Brixia

\item marted\`{\i} 11:15 - 13:30, laboratorio informatico 1 della sede Brixia
\end{stepitemize}
\end{stepitemize}

%

%TCIMACRO{\TeXButton{Transition: Box Out}{\transboxout}}%
%BeginExpansion
\transboxout
%EndExpansion%
%TCIMACRO{\TeXButton{EndFrame}{\end{frame}}}%
%BeginExpansion
\end{frame}%
%EndExpansion
%

%TCIMACRO{\TeXButton{BeginFrame}{\begin{frame}}}%
%BeginExpansion
\begin{frame}%
%EndExpansion
%

%TCIMACRO{\QTR{frametitle}{Cosa, chi, quando, dove, perch\'{e}, come}}%
%BeginExpansion
\frametitle{Cosa, chi, quando, dove, perch\'{e}, come}%
%EndExpansion
%

%TCIMACRO{\QTR{framesubtitle}{Cosa, chi, quando, dove}}%
%BeginExpansion
\framesubtitle{Cosa, chi, quando, dove}%
%EndExpansion
%

%TCIMACRO{\TeXButton{Transparency}{\setbeamercovered{transparent=20}}}%
%BeginExpansion
\setbeamercovered{transparent=20}%
%EndExpansion


\begin{stepitemize}
\item \textit{Ricevimento}

\begin{stepitemize}
\item prof. Paolo Falbo

gioved\`{\i}, 14:30 - 15:30

\begin{stepitemize}
\item Dipartimento di Economia e Management

via Porta Pile, 17/c, 3$%
%TCIMACRO{\U{b0}}%
%BeginExpansion
{{}^\circ}%
%EndExpansion
$ piano

25122 BRESCIA BS

\item via Microsoft Teams

\item e-mail: \textit{paolo.falbo@unibs.it}
\end{stepitemize}

\item prof. Cristian Pelizzari

luned\`{\i}, 16:45 - 17:45

\begin{stepitemize}
\item Dipartimento di Economia e Management

via Porta Pile, 17/c, 3$%
%TCIMACRO{\U{b0}}%
%BeginExpansion
{{}^\circ}%
%EndExpansion
$ piano

25122 BRESCIA BS

\item via Microsoft Teams

\item e-mail: c\textit{ristian.pelizzari@unibs.it}
\end{stepitemize}

\item \textit{dott. Niccol\`{o} Ghezzi}, e-mail:
\textit{niccolo.ghezzi@unibs.it}

sia in presenza che via Microsoft Teams \textbf{solo su appuntamento}
\end{stepitemize}
\end{stepitemize}

%

%TCIMACRO{\TeXButton{Transition: Box Out}{\transboxout}}%
%BeginExpansion
\transboxout
%EndExpansion%
%TCIMACRO{\TeXButton{EndFrame}{\end{frame}}}%
%BeginExpansion
\end{frame}%
%EndExpansion
%

%TCIMACRO{\TeXButton{BeginFrame}{\begin{frame}}}%
%BeginExpansion
\begin{frame}%
%EndExpansion
%

%TCIMACRO{\QTR{frametitle}{Perch\'{e} studiare MQF?}}%
%BeginExpansion
\frametitle{Perch\'{e} studiare MQF?}%
%EndExpansion
%

%TCIMACRO{\QTR{framesubtitle}{Dai modelli alle decisioni finanziarie}}%
%BeginExpansion
\framesubtitle{Dai modelli alle decisioni finanziarie}%
%EndExpansion
%

%TCIMACRO{\TeXButton{Transparency}{\setbeamercovered{transparent=20}}}%
%BeginExpansion
\setbeamercovered{transparent=20}%
%EndExpansion


\begin{stepitemize}
\item \textbf{Capire il rischio:} costruire scenari e misurare come
l'incertezza incide sui risultati finanziari.

\item \textbf{Prendere decisioni motivate:} confrontare alternative e vincoli
in casi concreti, dal rischio di credito alla gestione di attivit\`{a} e passivit\`{a}.

\item \textbf{Passare all'azione:} implementare gli algoritmi in Python, con
il supporto dell'IA, verificandone ipotesi e risultati.
\end{stepitemize}

%

%TCIMACRO{\TeXButton{Transition: Box Out}{\transboxout}}%
%BeginExpansion
\transboxout
%EndExpansion%
%TCIMACRO{\TeXButton{EndFrame}{\end{frame}}}%
%BeginExpansion
\end{frame}%
%EndExpansion
%

%TCIMACRO{\TeXButton{BeginFrame}{\begin{frame}}}%
%BeginExpansion
\begin{frame}%
%EndExpansion
%

%TCIMACRO{\QTR{frametitle}{Cosa, chi, quando, dove, perch\'{e}, come}}%
%BeginExpansion
\frametitle{Cosa, chi, quando, dove, perch\'{e}, come}%
%EndExpansion
%

%TCIMACRO{\QTR{framesubtitle}{Come}}%
%BeginExpansion
\framesubtitle{Come}%
%EndExpansion
%

%TCIMACRO{\TeXButton{Transparency}{\setbeamercovered{transparent=20}}}%
%BeginExpansion
\setbeamercovered{transparent=20}%
%EndExpansion


\begin{stepitemize}
\item \textit{Libro di testo}

\begin{stepitemize}
\item Manuale a cura dei docenti
\end{stepitemize}

\item \textit{Altro materiale}

\begin{stepitemize}
\item Slide delle presentazioni a lezione
\end{stepitemize}

\item \textit{Verifica dell'apprendimento}

\begin{stepitemize}
\item Per gli studenti che abbiano frequentato almeno il $75\%$ delle lezioni
la valutazione finale sar\`{a} determinata dalla media ponderata tra:

\begin{stepitemize}
\item soluzione dei cinque i casi applicativi assegnati durante le lezioni,
con peso $40\%$

\item prova orale finale sull'intero programma dell'insegnamento, con peso
$60\%$
\end{stepitemize}

\item In tutti gli altri casi, l'esame consiste in una prova orale
obbligatoria relativa all'intero programma
\end{stepitemize}
\end{stepitemize}

%

%TCIMACRO{\TeXButton{Transition: Box Out}{\transboxout}}%
%BeginExpansion
\transboxout
%EndExpansion%
%TCIMACRO{\TeXButton{EndFrame}{\end{frame}}}%
%BeginExpansion
\end{frame}%
%EndExpansion



\end{document}
